\documentclass[11pt,a4paper]{article}
\usepackage[margin=2.2cm]{geometry}
\usepackage{fontspec}
\setmainfont{Times New Roman}
\usepackage{amsmath,amssymb}
\usepackage{booktabs,tabularx,array}
\usepackage{enumitem}
\usepackage{xcolor}
\usepackage{float}
\usepackage{titlesec}
\usepackage[hidelinks]{hyperref}

\titleformat{\section}{\large\bfseries\color{black}}{\thesection}{0.6em}{}
\titleformat{\subsection}{\normalsize\bfseries\color{black}}{\thesubsection}{0.6em}{}

\setlength{\parindent}{0pt}
\setlength{\parskip}{4pt}
\linespread{1.2}
\renewcommand{\arraystretch}{1.15}
\renewcommand{\tabularxcolumn}[1]{>{\raggedright\arraybackslash}p{#1}}
\setlength{\emergencystretch}{2em}

\newcommand{\Cp}{\ensuremath{C_p}}
\newcommand{\CD}{\ensuremath{C_D}}
\newcommand{\ReD}{\ensuremath{\mathit{Re}}}

\pagestyle{plain}
\begin{document}

\begin{center}
{\LARGE\bfseries Cylinder Lab --- Operation Sheet}\\[3pt]
{\small What to do at the rig, in order \quad\textbar\quad RODH 115 (ex L50)
\quad\textbar\quad $d = 102$\,mm, $b = 460$\,mm, 28 tappings}\\[3pt]
{\footnotesize Monday 16 November 2026 \quad\textbar\quad Engineering Practice 2 -- Aerodynamics 2}
\end{center}

\vspace{2pt}

\section{Before you leave}

\begin{table}[H]
\centering
\caption{Kit and preparation}
\small
\begin{tabularx}{\textwidth}{@{}p{5.6cm} X c@{}}
\toprule
item & why & done \\
\midrule
prelab, printed or on the iPad & the supervisor may ask for the run matrix & $\square$ \\
\texttt{readings.csv} template on the laptop & 14 columns, 150 rows & $\square$ \\
\texttt{cylinder\_lab.m}, \texttt{cpUncertainty.m}, \texttt{runUncertainty.m} & all three in one folder &
$\square$ \\
MATLAB launched once at home (\texttt{cfg.demo = true}) & first start takes 40--60 s & $\square$ \\
Immutalog open, one notebook for this lab & an exported PDF cannot be edited & $\square$ \\
pen and a paper backup sheet & readings must survive a dead laptop & $\square$ \\
phone or calculator for the $\pm30$ check & one subtraction, one comparison & $\square$ \\
\bottomrule
\end{tabularx}
\end{table}

\section{Arrive, safety, switch-on}

The room door stays open whenever the tunnel is running --- a hard rule, not a preference. Nobody
reaches into the working section while the fan spins, and the fan comes to zero before anything is
mounted or touched.

Order: door open, fan to minimum, let it settle, level the board. Confirm all 28 tapping tubes plus
the two Pitot tubes are connected and none is pinched --- a single loose tube later appears as one
absurd \Cp{} and is easily mistaken for separation.

Record ambient temperature and barometric pressure before the first run, three readings each.

\section{Alignment using holes 1 and 7}

Holes 1 ($-30^\circ$) and 7 ($+30^\circ$) are mirrors about hole 2, so a correctly aligned cylinder
must give them the same reading. It is the only alignment check available, and a one-degree error
changes \CD{} by about 2\,\% --- the same size as the uncertainty you are reporting.

Procedure: run at the low-speed setting, read holes 1, 2 and 7 three times each, rotate by the
smallest increment the mount allows, repeat. Lock the position minimising $|h_1 - h_7|$.

Criterion, derived rather than guessed: with $B \approx 18.7$\,mm at low speed,
\begin{equation}
\Delta C_{p,30} = 2 g \Delta\theta \approx 2 \times 0.043 \times 1^\circ = 0.086
\quad\Rightarrow\quad |h_1 - h_7| \approx 0.086 \times 18.7 \approx 1.6\ \text{mm}
\end{equation}
so keeping the pair within about 1.5\,mm corresponds to less than one degree, the level the script
flags (\texttt{cfg.alignWarn = 1.0}). At high speed the same angular error gives four times the
column difference, so re-check once after raising the speed.

\section{Which Pitot tube is which}

The handout does not fix which of tubes 35 and 36 is static. Do not guess from the tube number,
and do not guess from the sign of the column --- decide physically. Hole 2 sits at $\theta = 0$,
so its surface pressure equals total pressure; the Pitot tube whose reading matches
$h(\text{hole } 2)$ is the total one.

At low speed read hole 2, then tubes 35 and 36. The tube within a fraction of a millimetre of hole 2
is total. The script repeats this test per run and prints a stagnation check line, so a flip between
runs shows up in the output --- that is a real fault, not a numerical curiosity.

\section{Run matrix}

Four runs are required, a fifth repeat is optional. Take the rpm values from Table 2 of the prelab.

\begin{table}[H]
\centering
\caption{Runs to record}
\small
\begin{tabular}{@{}c p{3.6cm} c c c p{5.0cm}@{}}
\toprule
run & boundary layer & $U_\infty$ & $q$ & column & purpose \\
\midrule
1 & smooth (natural) & 15.5 m/s & 145 Pa & 18.7 mm & attached, subcritical reference \\
2 & smooth (natural) & 30.5 m/s & 560 Pa & 72.4 mm & attached, higher \ReD \\
3 & tripped & 15.5 m/s & 145 Pa & 18.7 mm & post-critical drag drop \\
4 & tripped & 30.5 m/s & 560 Pa & 72.4 mm & post-critical, higher \ReD \\
5 (opt.) & as run 3 & 15.5 m/s & 145 Pa & 18.7 mm & repeat for a type A estimate \\
\bottomrule
\end{tabular}
\end{table}

Record the trip configuration in the logbook: type, position, method of attachment. Without it the
tripped runs cannot be compared with anyone else's, including your own next week.

\section{Reading routine per run}

One run is one pass over 28 tappings plus the two Pitot tubes. Keep the order identical every time
so a mistake appears as a shift rather than as noise.

\begin{enumerate}[leftmargin=1.5em,itemsep=3pt]
\item Wait for the speed to settle; check it has not drifted before starting and after finishing.
\item Read tubes 35 and 36 first, three repeats each. Everything downstream is scaled by this
difference, so it deserves the most care, not the least.
\item Read holes 1 to 28 in order, three repeats each. If a column oscillates, wait rather than
average a transient.
\item Read hole 2 again at the end. More than about 0.3\,mm between first and last reading means the
speed drifted during the pass; repeat the run.
\item Write the three repeats into the three adjacent columns of \texttt{readings.csv} while standing
at the rig. Do not carry numbers in your head to the laptop.
\end{enumerate}

Three repeats is the minimum that produces a type A term at all. A single reading forces the script
back to the conservative 0.5\,mm assumption and roughly triples the reported $u(\CD)$.

\section{Checks before leaving the rig}

All four take about five minutes with the tunnel still running. Each catches a failure that would
otherwise surface only in the analysis, when the rig is gone.

\begin{table}[H]
\centering
\caption{In-situ checks}
\small
\begin{tabularx}{\textwidth}{@{}p{3.6cm} X p{6.2cm}@{}}
\toprule
check & criterion & if it fails \\
\midrule
stagnation & $|h(\text{hole 2}) - h(\text{total})| \leq 0.3$\,mm & tubes swapped or leaking: re-identify, re-read \\
$\pm30$ symmetry & $|h_1 - h_7| \leq 1.5$\,mm at low speed & rotate the model, re-lock, redo the run \\
static drift & $|h(\text{hole 2})_{\text{start}} - h(\text{hole 2})_{\text{end}}| \leq 0.3$\,mm &
speed drifted: repeat the run \\
range & no column pinned at top or bottom & speed too high, or an air bubble: bleed the tube \\
\bottomrule
\end{tabularx}
\end{table}

If a run must be repeated, repeat it completely. Mixing half of one pass with half of another is the
most common way to produce a distribution that looks almost right.

\section{On the laptop, straight after the session}

Fill the wide template: one row per tapping, hole number in the \texttt{hole} column, the three
repeats in \texttt{r\_1}, \texttt{r\_2}, \texttt{r\_3}. Temperature and pressure occupy their own
three columns on the first row of each run; leaving them blank makes the script fall back to
21.0\,$^\circ$C and 101\,300\,Pa, which is not the same as writing the real values.

Set \texttt{cfg.demo = false} and \texttt{cfg.dataFile = 'readings.csv'}, then run the script.

\begin{table}[H]
\centering
\caption{What to look at in the console}
\small
\begin{tabularx}{\textwidth}{@{}p{3.4cm} p{4.2cm} X@{}}
\toprule
output & expected & meaning when outside the range \\
\midrule
stagnation check & within $\pm0.3$\,mm & tube identification wrong, or a leak \\
\texttt{dCp30} / ring rotation & $|\Delta\theta| < 1^\circ$ & model still misaligned; separation angle suspect \\
$u(\ReD')/\ReD'$ & 1--2\,\% low, 0.5\,\% high & far larger means the column was read to 1\,mm \\
$u(\CD')$ & 0.01--0.05, i.e.\ 1--5\,\% & a jump of 2--3$\times$ means repeats were not independent \\
$u(\theta_{\text{sep}})$ & 2--5$^\circ$ for attached runs & 20$^\circ$ or more: criterion failed, quote the peak angle \\
eq(2) versus integrated \CD & within about 10\,\% & a larger gap points at a bad tapping \\
\bottomrule
\end{tabularx}
\end{table}

Save the figures and the CSV, and export the logbook the same evening. The deadline is 23:00 the day
after the lab; an exported notebook can no longer be edited.

\section{What to record, and where it goes in the code}

The table below joins the two halves of the session: what you write at the rig, and what the script
reads afterwards. Column names are the ones in \texttt{readings.csv} and \textbf{must not be
changed} --- the file is read with \texttt{VariableNamingRule}, so a renamed column disappears.

\begin{table}[H]
\centering
\caption{Rig record $\rightarrow$ column of \texttt{readings.csv}}
\small
\begin{tabularx}{\textwidth}{@{}p{3.0cm} p{3.6cm} p{2.5cm} X@{}}
\toprule
item & column & how to write & note \\
\midrule
run number & \texttt{run} & 1--5 & all 30 rows of a run (28 holes + 2 tubes) share it \\
boundary layer & \texttt{tripped} & 0 or 1 & 0 smooth, 1 tripped; label only, no arithmetic \\
rpm & \texttt{rpm} & integer & a record, not an input: speed comes from the Pitot difference \\
scale division & \texttt{scale\_mm} & e.g.\ 1.0 & mm per division, sets $u_B=(s/2)/\sqrt{3}$; one row per run is enough \\
air temperature & \texttt{Tair\_C\_1..3} & 21.0 / 21.1 / 20.9 & three readings; one is also accepted \\
barometric pressure & \texttt{pBaro\_Pa\_1..3} & 101300 & three readings, in Pa, \textbf{not hPa} \\
hole & \texttt{hole} & 1--28 / 35 / 36 & 35 and 36 are the Pitot tubes; do not write 29, 30 \\
reading & \texttt{r\_1}, \texttt{r\_2}, \texttt{r\_3} & mm & the repeats; leave unread ones blank \\
\bottomrule
\end{tabularx}
\end{table}

\begin{table}[H]
\centering
\caption{\texttt{cfg} fields to check before the run}
\small
\begin{tabularx}{\textwidth}{@{}p{3.6cm} p{2.7cm} X@{}}
\toprule
field & default & when to change it \\
\midrule
\texttt{cfg.demo} & \texttt{true} & must be \texttt{false} in the lab, or the script reads synthetic data \\
\texttt{cfg.dataFile} & \texttt{'readings.csv'} & only if the file is renamed \\
\texttt{cfg.units} & \texttt{'mm'} & \texttt{'m'} only if the board is graduated in metres \\
\texttt{cfg.uRead} & 0.5 & applies only where a tapping has one reading; normally leave it \\
\texttt{cfg.uScale} & 1.0 & if the division is not 1\,mm and \texttt{scale\_mm} is blank \\
\texttt{cfg.uTheta} & 0.29 & only if the machining tolerance changes ($0.5^\circ/\sqrt{3}$) \\
\texttt{cfg.uTair}, \texttt{cfg.uPbaro} & 0.3, 50 & only if the thermometer or barometer changes \\
\texttt{cfg.alignWarn} & 1.0 & the alignment warning, in degrees; pairs with the 1.5\,mm criterion \\
\texttt{cfg.thetaSepManual} & \texttt{[]} & per-run override, e.g.\ \texttt{[80 88 NaN 124]}; \texttt{NaN} keeps the automatic value \\
\texttt{cfg.sepUncK} & 2 & 2 for a 95\,\% figure, 1 for a bare $1\sigma$ \\
\texttt{cfg.figFormat} & \texttt{'pdf'} & \texttt{'png'} only if a raster is wanted \\
\texttt{cfg.uSigma} & \texttt{[1 2]} & add 3 for a three-level bar \\
\texttt{cfg.pitotOrder} & \texttt{'auto'} & force \texttt{'tube35static'} / \texttt{'tube36static'} if the auto test disagrees with your observation \\
\texttt{cfg.TairDefault} & 21.0 & used when the temperature column is blank --- \textbf{set it to the real value} \\
\texttt{cfg.pBaroDefault} & 101300 & same, in Pa \\
\texttt{cfg.nFine} & 0.25 & refined integration grid in degrees; leave it \\
\texttt{cfg.outDir} & \texttt{results} & output directory \\
\bottomrule
\end{tabularx}
\end{table}

The physical constants (\texttt{dia}, \texttt{tunnelB}, \texttt{rhoEth}, \texttt{g}, \texttt{Rair},
\texttt{mu0}, \texttt{T0}, \texttt{Suth}) come from the handout and must not be touched; the same
holds for the array of tapping angles.

\begin{table}[H]
\centering
\caption{What a blank or a missing row falls back to}
\small
\begin{tabularx}{\textwidth}{@{}p{4.4cm} p{4.2cm} X@{}}
\toprule
situation & consequence & say this in the report \\
\midrule
all \texttt{Tair\_C\_*} blank & \texttt{TairDefault} = 21.0\,$^\circ$C is used & $\rho$ rests on a default temperature; state it \\
all \texttt{pBaro\_Pa\_*} blank & \texttt{pBaroDefault} (101300\,Pa) is used & same \\
\texttt{scale\_mm} blank & \texttt{cfg.uScale} = 1.0\,mm is used & $u_B$ becomes 0.289\,mm \\
one repeat only & that point falls back to 0.5\,mm; no type A term & $u(\CD)$ roughly triples \\
a tapping with no reading & the run prints \texttt{[skip]} and is dropped & fill it in and re-run; never use 0 as a placeholder \\
35 or 36 missing & the run prints \texttt{[skip]} & same \\
a whole row missing & same & same \\
\bottomrule
\end{tabularx}
\end{table}

One sentence: \textbf{the CSV carries the data, \texttt{cfg} carries the model}. Readings go into the
CSV, constants and thresholds into \texttt{cfg}, and the two should not fight each other. The result
CSV has 29 columns, but only seven numbers belong in the logbook: $\CD'$, $\CD$ from eq.(2),
$\CD$ from the integral, $V/V'$, $\theta_{\text{sep}}$, \ReD{}', plus $u(\CD')$,
$u(\theta_{\text{sep}})$ and $\Delta\theta$.

\section{Troubleshooting}

\begin{table}[H]
\centering
\caption{Common failures}
\small
\begin{tabularx}{\textwidth}{@{}p{5.2cm} X p{5.4cm}@{}}
\toprule
symptom & most likely cause & action \\
\midrule
one \Cp{} far outside its neighbours & air bubble or kinked tube & bleed and re-read that tapping only \\
all \Cp{} shifted the same way & static tube wrong, or board re-levelled mid-run & re-zero, repeat the pass \\
\Cp{} at hole 2 below the total tube & hole 2 is not at stagnation & check alignment, re-identify the tubes \\
separation near $180^\circ$ on a tripped run & criterion found no flat region &
expected at high blockage; report the peak angle and say so \\
\ReD{} outside $10^5$--$10^6$ & speed not what the rpm column claims & re-read the dynamic head, correct the entry \\
script stops with a missing-module note & one uncertainty file absent & copy both files next to
\texttt{cylinder\_lab.m}; the run still completes without them \\
\bottomrule
\end{tabularx}
\end{table}

\section{Timing}

\begin{table}[H]
\centering
\caption{Realistic budget for a two-hour session}
\small
\begin{tabularx}{\textwidth}{@{}X c p{7.4cm}@{}}
\toprule
block & minutes & note \\
\midrule
safety, switch-on, levelling & 10 & door open, fan at minimum \\
alignment with holes 1 and 7 & 15 & iterate until $|h_1 - h_7| \leq 1.5$\,mm \\
Pitot identification & 5 & hole 2 versus 35 and 36 \\
run 1 (smooth, low) & 12 & 30 readings, three repeats each \\
run 2 (smooth, high) & 12 & same pass, higher speed \\
fit or remove the trip & 10 & record type and position \\
run 3 (tripped, low) & 12 & expect a visibly different plateau \\
run 4 (tripped, high) & 12 & same \\
optional run 5 & 12 & repeat of run 3 \\
final checks and pack down & 10 & four checks, fan to zero, tubes off \\
\bottomrule
\end{tabularx}
\end{table}

\section{One-page checklist}

\begin{table}[H]
\centering
\caption{Sequence at the rig}
\small
\begin{tabularx}{\textwidth}{@{}c X p{5.0cm} c@{}}
\toprule
\# & step & done when & $\checkmark$ \\
\midrule
1 & door open, fan at minimum, board levelled & all 30 tubes connected, none pinched & $\square$ \\
2 & $T$ and $p$ recorded, three repeats each & values in the first row of the run & $\square$ \\
3 & model aligned using holes 1 and 7 & $|h_1 - h_7| \leq 1.5$\,mm at low speed & $\square$ \\
4 & Pitot tubes identified & $|h(\text{hole 2}) - h(\text{total})| \leq 0.3$\,mm & $\square$ \\
5 & run 1, smooth low & hole 2 start and end within 0.3\,mm & $\square$ \\
6 & run 2, smooth high & same & $\square$ \\
7 & trip fitted, configuration written & logbook entry complete & $\square$ \\
8 & run 3, tripped low & plateau visible in the tail & $\square$ \\
9 & run 4, tripped high & same & $\square$ \\
10 & optional run 5 & repeat of run 3 & $\square$ \\
11 & four in-situ checks & all four pass & $\square$ \\
12 & pack down & fan to zero, tubes off, section clear & $\square$ \\
13 & \texttt{readings.csv} filled, script run & no warnings in the console & $\square$ \\
14 & figures and CSV saved, logbook exported & PDF before 23:00 the next day & $\square$ \\
\bottomrule
\end{tabularx}
\end{table}

\end{document}
